\section{Truyền tin cậy trên UDP}
\label{sec:tin-cay}

UDP chỉ ``cố gắng gửi'' (\emph{best effort}): gói có thể bị mất, tới sai thứ tự hoặc bị lặp, và bên gửi cũng
không biết mạng còn chịu được bao nhiêu dữ liệu. Chương này trình bày cách tầng \fname{internal/transport} của
nhóm biến UDP thành một kênh truyền \textbf{tin cậy, đúng thứ tự trong từng stream, có mã hóa, không bị HoL và tự
điều chỉnh tốc độ}. Nhóm thiết kế dựa theo QUIC \cite{rfc9000,rfc9002} nhưng lược bớt cho vừa phạm vi bài tập
lớn. \Cref{tab:van-de-udp} tóm tắt các vấn đề của UDP và cơ chế nhóm dùng để giải quyết, các mục sau sẽ đi vào
chi tiết từng cơ chế.

\begin{table}[H]
  \centering
  \caption{Các vấn đề của UDP và cơ chế giải quyết}
  \label{tab:van-de-udp}
  \begin{tabularx}{\linewidth}{lY}
    \toprule
    \textbf{Vấn đề} & \textbf{Cơ chế} \\
    \midrule
    Mất gói & SACK; phát hiện mất theo số đếm, theo thời gian, theo RTO; chỉ gửi lại các frame của gói mất,
    trong gói mới với \texttt{packet\_num} mới (\cref{subsec:mat-goi}) \\
    Sai thứ tự & Ghép theo \texttt{offset}, riêng từng stream (\cref{subsec:ghep}) \\
    Lặp gói & Bỏ \texttt{packet\_num} đã thấy; bỏ đoạn có offset đã giao \\
    Gửi quá nhanh & Chỉ gửi khi $\mathit{inflight} + 1\ \text{gói} \le \min(\mathit{cwnd},
    \mathit{window})$; hàng đợi đầy thì chặn ứng dụng (\cref{subsec:dieu-toc}, \cref{sec:hol}) \\
    Gói rác, gói giả mạo & Sai magic/version/\texttt{conn\_id} hoặc giải mã GCM thất bại thì bỏ im lặng
    (\cref{sec:bao-mat}) \\
    Bên kia chết & PING sau 2 s im lặng; không nhận được gói hợp lệ nào trong 10 s thì đóng kết nối
    (\cref{subsec:song-chet}) \\
    \bottomrule
  \end{tabularx}
\end{table}

\subsection{Giao diện cho tầng ứng dụng}
\label{sec:api}

\begin{lstlisting}[caption={Interface tầng transport cung cấp cho tầng ứng dụng (\fname{api.go})},
  label=lst:endpoint]
type Endpoint interface {
	SendMessage(b []byte) error            // 1 message = 1 high-priority stream
	OpenStream(urgent bool) io.WriteCloser // long stream (e.g. a file); Close() sends FIN
	Recv() <-chan Delivery                 // in-order chunks, per stream
	Done() <-chan struct{}                 // closed on disconnect / idle timeout / Close()
	Stats() Stats                          // RTT, cwnd, in-flight bytes, lost packets...
	Close() error
}

type Delivery struct {
	StreamID uint32
	Data     []byte
	Fin      bool // last chunk of the stream
}
\end{lstlisting}

\Cref{lst:endpoint} là interface mà tầng transport cung cấp cho tầng ứng dụng. Dữ liệu nhận được trả về dưới
dạng các \fname{Delivery}. Các đoạn thuộc cùng một stream luôn đúng thứ tự, còn các đoạn của những stream khác
nhau thì có thể xen kẽ nhau, stream nào đủ dữ liệu trước thì được giao trước. Với interface này, việc gửi một
file chỉ cần một dòng \fname{io.Copy(ep.OpenStream(false), file)}.

Ngoài bản chạy trên UDP thật (\fname{*Conn}), nhóm viết thêm \fname{FakeEndpoint} nối hai đầu ngay trong bộ
nhớ. Nhờ có bản giả này, phần tầng ứng dụng có thể được viết và test song song trong lúc tầng transport chưa làm
xong. Hàm \fname{Stats()} trả về các số liệu mạng như RTT, cwnd, số gói mất để giao diện hiển thị trạng thái kết
nối.

\subsection{Mô hình concurrency}
\label{subsec:concurrency}

\begin{figure}[H]
  \centering
  \begin{tikzpicture}[
      font=\small,
      box/.style={draw, rounded corners, fill=blue!6, align=center, minimum height=1cm, text width=2.3cm},
      run/.style={draw, very thick, rounded corners, fill=orange!12, align=center, minimum height=2.8cm,
        text width=3.2cm},
      ch/.style={-{Stealth}, thick},
      lab/.style={font=\scriptsize, midway, above, align=center}]
    \node[run] (run) at (0,0) {goroutine \texttt{run()}\\[3pt]\scriptsize sở hữu toàn bộ state:
      unacked, hàng đợi, cwnd, RTT, stream};
    \node[box] (app)  at (-5.5, 0.85) {Ứng dụng\\\scriptsize \texttt{SendMessage}, \texttt{Write}};
    \node[box] (appr) at (-5.5,-0.85) {Ứng dụng\\\scriptsize đọc \texttt{Recv()}};
    \node[box] (rl)   at ( 5.5, 0.85) {\texttt{readLoop}\\\scriptsize đọc socket, tra \texttt{conn\_id}};
    \node[box] (sock) at ( 5.5,-0.85) {socket UDP\\\scriptsize \texttt{WriteToUDP}};
    \node[box] (tick) at (0,-2.4) {ticker 10/100 ms};
    \draw[ch] (app.east) -- (app.east -| run.west) node[lab] {urgentCh\\bulkCh};
    \draw[ch] (appr.east -| run.west) -- (appr.east) node[lab] {recvCh};
    \draw[ch] (rl.west) -- (rl.west -| run.east) node[lab] {in};
    \draw[ch] (sock.west -| run.east) -- (sock.west) node[lab] {gói đã\\mã hóa};
    \draw[ch] (tick) -- (run);
  \end{tikzpicture}
  \caption{Mỗi kết nối có đúng một goroutine sở hữu state; goroutine khác chỉ giao tiếp qua channel}
  \label{fig:concurrency}
\end{figure}

Một vấn đề nhóm phải tính tới sớm là nhiều goroutine cùng truy cập state của một kết nối: ứng dụng gọi
\fname{Write} để gửi, trong khi \fname{readLoop} đẩy gói nhận được vào. Nếu dùng mutex để bảo vệ thì rất dễ quên
khóa hoặc bị deadlock. Vì vậy nhóm chọn cách cho mỗi kết nối (\fname{Conn}) có \textbf{đúng một goroutine}
\fname{run()} giữ toàn bộ state, bao gồm: các gói chưa được ACK, các hàng đợi gửi, cwnd, giá trị ước lượng RTT
và các stream đang nhận. Các goroutine khác chỉ được giao tiếp với nó qua channel (\cref{fig:concurrency}).

Bên trong, \fname{run()} là một vòng lặp \fname{select} chờ các sự kiện: có gói tới, ứng dụng đưa frame mới
xuống, ứng dụng sẵn sàng nhận thêm dữ liệu, tới nhịp đồng hồ, hoặc có yêu cầu đóng kết nối. Sau mỗi sự kiện,
\fname{run()} gọi \fname{flush()} để gửi đi những gì đang được phép gửi.

Nhờ thiết kế này, phần phức tạp nhất là truyền tin cậy và điều khiển tắc nghẽn \textbf{không cần dùng mutex} và
không thể xảy ra data race. Nhóm đã kiểm tra lại bằng \fname{go test -race}.

Phía server dùng chung một socket UDP cho tất cả client. \fname{Mux} có một goroutine \fname{readLoop} đọc
socket, xem \texttt{conn\_id} trong header (phần này không mã hóa) và chuyển gói vào channel của đúng
\fname{Conn}. Như vậy 20 client sẽ có 20 goroutine \fname{run()} chạy độc lập, kết nối này chậm cũng không làm
chậm các kết nối khác.

Nhịp đồng hồ dùng để phát hiện mất gói theo thời gian và để gửi PING. Khi có gói đang chờ ACK, đồng hồ chạy mỗi
10 ms; khi rảnh thì chỉ chạy mỗi 100 ms để đỡ tốn CPU (thử thách~\ref{tt:cpu} ở \cref{sec:thu-thach}).

\subsection{Đường đi của dữ liệu}
\label{subsec:duong-di}

Mục này mô tả dữ liệu đi qua tầng transport như thế nào, từ lúc ứng dụng gọi hàm gửi cho tới khi bên kia nhận
được.

\paragraph{Phía gửi.}
\begin{enumerate}
  \item \fname{SendStream.Write} cắt dữ liệu thành các frame không quá 1000 byte, gắn \texttt{stream\_id} và
        \texttt{offset}, rồi đẩy vào \texttt{urgentCh} nếu là tin nhắn hoặc \texttt{bulkCh} nếu là file.
  \item \fname{run()} chuyển frame vào hàng đợi tương ứng. Khi hàng đợi file đã đầy 256 frame, \fname{run()}
        tạm ngừng nhận từ \texttt{bulkCh}, nên lệnh \fname{Write} của ứng dụng bị chặn lại. Cơ chế này gọi là
        \emph{backpressure}, nhờ nó mà gửi một file 200 MB cũng không làm bộ nhớ phình to.
  \item \fname{flush()}: nếu cửa sổ còn chỗ thì lấy frame theo thứ tự ưu tiên \texttt{retransQ} (cần gửi lại)
        $\to$ \texttt{urgentQ} (tin nhắn) $\to$ \texttt{bulkQ} (file) cho tới khi đầy gói (\cref{sec:hol}).
  \item \fname{sendPacket()}: ghép ACK đang cần gửi (nếu có) với các frame, gán một \texttt{packet\_num} mới,
        mã hóa rồi gửi bằng \fname{WriteToUDP}. Gói được lưu vào \texttt{unacked[pn]} cùng với thời điểm gửi và
        danh sách frame, để sau này biết cần gửi lại những gì nếu gói bị mất.
\end{enumerate}

\paragraph{Phía nhận.} Gói đi từ \fname{readLoop} vào \fname{handlePacket}. Hàm này kiểm tra header, giải mã
(header được dùng làm AAD, xem \cref{sec:bao-mat}), tách các frame và ghi nhận \texttt{packet\_num} để ACK sau.
Mỗi loại frame được xử lý riêng:
\begin{itemize}
  \item Frame ACK đi vào \fname{onAck}: xóa các gói đã được xác nhận khỏi \texttt{unacked}, đo RTT, tăng cwnd
        và kiểm tra xem có gói nào bị mất không.
  \item Frame STREAM đi vào \fname{recvStream.push} để ghép theo offset, rồi được giao lên cho ứng dụng.
\end{itemize}

\subsection{ACK và phát hiện mất gói}
\label{subsec:mat-goi}

Bên nhận gửi ACK ngay trong lần \fname{flush()} kế tiếp chứ không trì hoãn như TCP. Nếu lúc đó có dữ liệu
cần gửi thì ACK được gửi kèm luôn trong gói dữ liệu. Ngoài ra, sau khi xử lý một gói, \fname{run()} xử lý luôn
tối đa 32 gói đang chờ trong channel rồi mới flush, nên một ACK có thể xác nhận được nhiều gói cùng lúc. Mỗi ACK
chứa tối đa 8 khoảng, đủ để mô tả các chỗ bị thiếu gần nhất.

\paragraph{Ước lượng RTT.} RTT (\emph{round-trip time}) là thời gian từ lúc gửi một gói tới lúc nhận được ACK
của nó. Mỗi khi nhận ACK, nhóm lấy thời gian của gói mới nhất được xác nhận làm một mẫu RTT $R$, rồi cập nhật
giá trị trung bình $\mathit{SRTT}$ và độ dao động $\mathit{RTTVAR}$ theo RFC~6298 \cite{rfc6298}:
\begin{equation}
  \mathit{RTTVAR} \leftarrow \tfrac34 \mathit{RTTVAR} + \tfrac14 \lvert \mathit{SRTT} - R \rvert, \qquad
  \mathit{SRTT} \leftarrow \tfrac78 \mathit{SRTT} + \tfrac18 R
\end{equation}

\paragraph{Ba cách kết luận một gói đã mất.} Theo RFC~9002 \cite{rfc9002}, nhóm dùng ba cách sau để quyết định
một gói đã bị mất:
\begin{enumerate}
  \item \textbf{Theo số đếm:} gói $pn$ chưa được ACK trong khi gói $pn+3$ hoặc lớn hơn đã được ACK, tức là gói
        $pn$ đã bị ``vượt mặt'' 3 gói. Nhóm chọn ngưỡng 3 để chừa chỗ cho trường hợp mạng đảo thứ tự gói nhẹ.
  \item \textbf{Theo thời gian:} đã có gói gửi sau nó được ACK, và nó đã chờ lâu hơn
        $\tfrac98 \max(\mathit{SRTT}, R)$. Cách này cần thiết khi cwnd chỉ còn 2--3 gói, vì khi đó sẽ không bao
        giờ có đủ 3 gói để vượt mặt (thử thách~\ref{tt:cubic-cham}).
  \item \textbf{Theo RTO} (\emph{retransmission timeout}): không nhận được ACK nào trong khoảng thời gian
        \begin{equation}
          \mathit{RTO} = \mathit{SRTT} + \max(4\cdot\mathit{RTTVAR},\ 1\,\text{ms}) + 10\,\text{ms}
          \label{eq:rto}
        \end{equation}
        Đây là công thức PTO của QUIC. Khi chưa có mẫu RTT nào, RTO ban đầu được đặt là 500 ms.
\end{enumerate}

Khi một gói bị coi là mất, \textbf{chỉ các frame STREAM của nó} được đưa vào \texttt{retransQ} để gửi lại; frame
ACK và PING thì không cần gửi lại. Các frame này giữ đúng thứ tự gửi ban đầu và giữ nguyên offset, nhưng được
đóng vào một gói mới với \texttt{packet\_num} mới. \Cref{fig:seq-loss} minh họa toàn bộ quá trình này.

\begin{figure}[H]
  \centering
  \begin{tikzpicture}[font=\small, x=0.9cm, y=0.5cm,
      msg/.style={-{Stealth}, thick}, lost/.style={-{Stealth}, thick, dashed, red!70!black},
      ack/.style={-{Stealth}, thick, blue!60!black},
      lab/.style={font=\scriptsize, above, sloped, pos=0.5},
      note/.style={font=\scriptsize, align=left, fill=yellow!15, inner sep=3pt}]
    \node[draw, fill=blue!6, minimum width=2cm] at (0,0) {Bên gửi};
    \node[draw, fill=blue!6, minimum width=2cm] at (8,0) {Bên nhận};
    \draw[gray] (0,-0.8) -- (0,-14.5);
    \draw[gray] (8,-0.8) -- (8,-14.5);
    \draw[lost] (0,-1.5) -- (4,-2.5) node[lab] {\#41 STREAM(id=5, off=0)};
    \node[red!70!black, font=\large] at (4.1,-2.55) {$\times$};
    \draw[msg] (0,-3) -- (8,-4.5) node[lab] {\#42 STREAM(id=5, off=1000)};
    \draw[msg] (0,-4.2) -- (8,-5.7) node[lab, pos=0.15] {\#43};
    \draw[msg] (0,-5.4) -- (8,-6.9) node[lab, pos=0.15] {\#44};
    \node[note, right] at (8.15,-4.5) {stream 5: giữ off=1000\\ở \texttt{pending}, chờ off=0};
    \draw[ack] (8,-7.3) -- (0,-8.8) node[lab] {ACK [42..44]};
    \node[note, left] at (-0.15,-8.8) {\#41 bị 3 gói\\vượt mặt $\Rightarrow$ mất};
    \draw[msg] (0,-9.8) -- (8,-11.3) node[lab] {\#45 STREAM(id=5, off=0): số gói mới};
    \node[note, right, fill=green!12] at (8.15,-11.3) {lấp chỗ trống $\Rightarrow$\\giao off=0, off=1000};
    \draw[ack] (8,-11.9) -- (0,-13.4) node[lab] {ACK [42..45]};
    \node[note, left] at (-0.15,-13.4) {mẫu RTT của \#45\\không lẫn với \#41};
  \end{tikzpicture}
  \caption{Một gói bị mất và được gửi lại (phát hiện theo số đếm)}
  \label{fig:seq-loss}
\end{figure}

\paragraph{RTO backoff.} Nếu hết RTO nhiều lần liên tiếp thì RTO được nhân đôi sau mỗi lần, tối đa 5 s; khi có
ACK mới thì RTO trở lại bình thường. Mục đích là khi mạng chết hẳn, bên gửi không gửi gói dồn dập làm mạng càng
tệ hơn. Còn nếu mất gói được phát hiện theo số đếm hoặc theo thời gian thì nghĩa là ACK vẫn về, mạng vẫn còn
hoạt động, nên trường hợp này không tăng backoff.

\paragraph{Vì sao không dùng RTO tối thiểu 200 ms như TCP?} TCP cần một mức sàn lớn cho RTO vì bên nhận TCP được
phép trì hoãn ACK tới 200--500 ms. Bên nhận trong giao thức của nhóm thì ACK ngay, nên nhóm dùng
công thức~\eqref{eq:rto} mà không đặt sàn 200 ms.

Nhóm đã đo thử để so sánh. Trên loopback với tỉ lệ mất gói 10\%, nếu đặt sàn 200 ms thì throughput giảm khoảng
6 lần (141 KB/s so với 826 KB/s). Nguyên nhân là gần như mọi lần mất gói ở cuối cửa sổ đều phải chờ đủ 200 ms,
trong khi RTT thực tế chỉ khoảng 0,4 ms.

Ngược lại, nhóm cũng thử bỏ hẳn khoảng đệm, tức là dùng $\mathit{RTO} = \mathit{SRTT} + 4\cdot\mathit{RTTVAR}$.
Khi đường truyền rất ổn định, RTTVAR giảm dần về 0 nên RTO gần bằng RTT, và chỉ cần một ACK về chậm một chút là
gói bị tính là mất oan. Trên nút cổ chai 2 MB/s, lỗi này kéo throughput xuống còn 0,87 MB/s. Sau khi thêm khoảng
đệm 1 ms và 10 ms như trong công thức~\eqref{eq:rto}, throughput đạt 2,01 MB/s.

\subsection{Ghép stream theo offset}
\label{subsec:ghep}
Ở phía nhận, mỗi stream có một \fname{recvStream} riêng gồm hai thành phần:
\begin{itemize}
  \item \texttt{next}: offset tiếp theo cần giao cho ứng dụng; mọi byte trước vị trí này đều đã được giao.
  \item \texttt{pending}: các đoạn dữ liệu tới sớm, được sắp theo offset.
\end{itemize}
Khi nhận một frame chứa các byte trong khoảng $[\mathit{off}, \mathit{end})$, có ba trường hợp:

\begin{itemize}
  \item $\mathit{end} \le \mathit{next}$: dữ liệu này đã được giao rồi (do gửi lại hoặc gói bị lặp), nên bỏ đi.
  \item $\mathit{off} > \mathit{next}$: dữ liệu tới sớm, phía trước vẫn còn chỗ trống, nên cất tạm vào
        \texttt{pending}.
  \item Các trường hợp còn lại: frame nối tiếp được \texttt{next} (có thể chồng lên một phần). Khi đó nhóm cắt bỏ
        phần đã giao, giao phần mới cho ứng dụng, rồi giao tiếp các đoạn trong \texttt{pending} vừa trở nên liền
        mạch.
\end{itemize}

Stream được coi là hoàn tất khi đã giao đủ mọi byte cho tới FIN. Frame mang cờ FIN có thể tới trước phần dữ liệu
phía trước nó; trong trường hợp này nhóm ghi nhận độ dài của stream rồi chờ cho tới khi nhận đủ.

\subsection{Điều khiển luồng}
\label{subsec:dieu-toc}
Điều khiển luồng nhằm đảm bảo bên gửi không gửi nhanh hơn khả năng xử lý của bên nhận. Mỗi bên nhận có một bộ
đệm 1 MB, chứa dữ liệu đã ghép xong nhưng ứng dụng chưa đọc và dữ liệu tới sớm trong \texttt{pending}. Mỗi ACK
mang trường \texttt{window} là số byte còn trống trong bộ đệm này. Bên gửi chỉ gửi thêm một gói khi
\begin{equation}
  \min(\mathit{cwnd}, \mathit{window}) - \mathit{bytesInFlight} \ge 1200
\end{equation}
Như vậy, nếu bên nhận đọc chậm (ví dụ trình duyệt bị treo) thì bên gửi sẽ tự chậm lại theo, thay vì gửi tới mức
làm tràn bộ nhớ của bên nhận.

Có một trường hợp nhóm phải xử lý thêm: khi cửa sổ đã báo cho bên gửi xuống dưới 1/4 nhưng sau đó bộ đệm trống lại
trên 1/2, bên nhận chủ động gửi một ACK để cập nhật cửa sổ, nếu không bên gửi sẽ đứng chờ không cần thiết. Ngoài
ra, một ACK cũ tới sau ACK mới thì không được ghi đè giá trị cửa sổ.

\subsection{Phát hiện mất kết nối}
\label{subsec:song-chet}
UDP không có khái niệm kết nối, nên nhóm phải tự xác định khi nào bên kia đã rời đi hoặc mất mạng. Nhóm dùng các
cơ chế sau:
\begin{itemize}
  \item \textbf{Keep-alive:} nếu không gửi gì trong 2 s thì gửi một gói PING. Vì PING cần được ACK nên hai bên
        luôn biết bên kia còn hoạt động.
  \item \textbf{Idle timeout:} nếu không nhận được gói hợp lệ nào trong 10 s thì coi như bên kia đã chết. Khi đó
        kết nối bị đóng và ứng dụng được báo qua \fname{Done()}.
  \item \textbf{CLOSE:} khi chủ động đóng kết nối, bên đóng gửi CLOSE để bên kia biết ngay thay vì phải đợi
        10 s. CLOSE chỉ gửi một lần, không gửi lại; nếu gói này bị mất thì bên kia vẫn tự đóng sau idle timeout.
  \item \textbf{PING đầu tiên:} server chỉ biết địa chỉ UDP của client khi nhận được gói đầu tiên từ client. Vì
        vậy client gửi lại PING đầu sau 200, 400, 800 ms\dots{} (tối đa 2 s) cho tới khi nhận được phản hồi từ
        server (thử thách~\ref{tt:ping-dau}).
  \item \textbf{Gói hỏng:} gói sai magic hoặc version, có \texttt{conn\_id} lạ, frame bị cụt, type không hợp lệ
        hoặc giải mã thất bại đều bị bỏ đi. Chương trình không bị panic và kết nối không bị đóng, nhờ vậy kẻ gửi
        gói rác không thể làm ngắt kết nối của người khác.
\end{itemize}
